% Generated by analysis/scripts/export_ten_cases.py
\subsection*{统一结果卡}
\begingroup\small
\noindent\begin{tabularx}{\textwidth}{@{}p{24mm}X@{}}\toprule
技术/模式 & SRAM ACIM；二进制9T1C电荷域；128行求和；8位平面并行\\
逻辑粒度/聚合 & B\_S=128 Byte；B\_R=16 Byte。同一输入行的16个输出权重，对齐16 Byte；1024事务覆盖整矩阵\\
资源/相对R0 & 128 SAR、16条两拍重构通道、128对真实写驱动；沿用R0求值组织，前端合并预算\\
服务口径 & 无周期维护；普通SRAM需持续供电\\
主导因素/选择 & 参考20 ns合并前端与20 ns SAR各占读服务39.9\%；普通写5 ns完整周期与DCIM同源，不再加控制拍\\
关键对照 & 固定外围扫前端10/20/50 ns；64写驱动对照单列\\
\bottomrule\end{tabularx}
\vspace{3mm}
\begin{center}\begin{tabular}{@{}lrrr@{}}\toprule
成对条件情景 & $\rho$ (MB/s) & $\tau$ (MB/s) & $\mathrm{RI}^{*}$\\\midrule
短预算 & 83.12 & 7273 & 0.01143\\
参考（推荐） & 39.88 & 3200 & 0.01246\\
长预算 & 16.62 & 1600 & 0.01039\\
\bottomrule\end{tabular}\end{center}
\noindent 参考原始占用：streaming 3210 ns；resident 5 ns。
\par\noindent 范围为有限成对条件情景极值，不是物理不确定性包络。1 MB=$10^6$ Byte；整矩阵16384 Byte=16 KiB。源定位见本例证据笔记；公共基线shared\_baseline。
\endgroup
\clearpage
